\documentclass[10pt,a4paper]{amsart}
\usepackage[margin=19mm]{geometry}
\usepackage{amsmath,amssymb,mathtools}
\usepackage[scr=rsfs,cal=euler]{mathalfa}
\usepackage{xfrac}
\usepackage[unicode,hidelinks,bookmarks=false]{hyperref}
\newcommand{\ie}{i.e.\ }
\newcommand{\cf}{cf.\ }
\newcommand{\resp}{respectively\ }
\newcommand{\Subsection}{\subsection}
\providecommand{\nobreakdash}{\nobreak\hbox{-}}
\setlength{\emergencystretch}{2em}
\multlinegap0pt
\usepackage{amsthm}
\usepackage{bm}
\usepackage{enumitem}
\theoremstyle{plain}
\newtheorem{defn}{Definition}
\newtheorem{lemma}{Lemma}
\newtheorem{example}{Example}
\newtheorem{theorem}{Theorem}
\newtheorem{remark}{Remark}
\newtheorem{corollary}{Corollary}
\newtheorem{note}{Note}
\newcommand{\floor}[1]{\left\lfloor #1 \right\rfloor}

\newcommand\at[2]{\left.#1\right|_{#2}}
\newlist{mycases}{enumerate}{1}
\setlist[mycases]{label=\textit{Case~\Roman*:},align=left,itemindent=0pt,leftmargin=0pt,labelwidth=-4pt}


\newlist{mysubcases}{enumerate}{1}
\setlist[mysubcases]{label=\textit{Sub-case~\arabic*:},align=left,itemindent=0pt,leftmargin=0.6cm,labelwidth=-4pt}
\title[Construction of 2-D Z-Complementary Array Code Sets with Flexible Lengths for Different System Requirements]{Construction of 2-D Z-Complementary Array Code Sets with Flexible Lengths for Different System Requirements: the 2-D ZCACS construction (Theorem 1)}
\author{Abhishek Roy}
\date{Source: arXiv:2109.00970v4 (CC BY 4.0)}
\begin{document}
\maketitle
\noindent\emph{Source and licence.} Construction of 2-D Z-Complementary Array Code Sets with Flexible Lengths for Different System Requirements, arXiv:2109.00970v4, distributed under the Creative Commons Attribution 4.0 International licence, http://creativecommons.org/licenses/by/4.0/, as recorded in the arXiv OAI metadata for this exact version and shown on the abstract page https://arxiv.org/abs/2109.00970v4 (v4 of 2 September 2025). Full rights notice: licenses/arxiv-2109.00970-CC-BY-4.0.txt.\par\smallskip
\noindent\textit{Editorial scope.} Counted unit: the paper\textquotesingle s main construction theorem, source label \textbackslash{}label\{thm\_ZCACS\} (source0.tex lines 703--705), printed as Theorem 1 of the article (version PDF page 15), delivered with its complete original proof, the whole of the author\textquotesingle s \textbackslash{}begin\{proof\} block at source0.tex lines 707--891 (the trivial-case count, the splitting of the inner sum into S\_1, S\_2, S\_3 and then T\_1, T\_2, the two case analyses on tau\_2 and tau\_1, and the concluding vanishing statement), as the single primary proof body.\par
\noindent\textit{Required context retained uncounted.} One source span: the statement of the author\textquotesingle s own IGC construction lemma (source0.tex lines 444--446, source label \textbackslash{}label\{Thm\_IGC\_gen\}), printed as Lemma 5 of the article, which the counted proof invokes twice; it is the only result of the author\textquotesingle s own development that the proof names. One disclosed adaptation of a cross-reference: the phrase using Lemma \textbackslash{}ref\{lemma\_mate\} at source line 868 points at the imported complementary-mate lemma (source label lemma\_mate, printed as Lemma 9, attributed by the author to Rathinakumar and Chaturvedi), which is not part of the author\textquotesingle s own development and is not reproduced here; that single command is delivered as \textbackslash{}href\{https://arxiv.org/abs/2109.00970v4\}\{9\}, a hyperlink to the version of record whose visible text is the lemma number the version PDF prints at that place. The two occurrences of \textbackslash{}ref\{Thm\_IGC\_gen\} resolve normally against the delivered context lemma, whose printed number is restored to 5. These are the only replacements in the package; they are recorded as a reference-only adaptation in the delivery evidence, and every other symbol of those proof lines is unchanged.\par
\noindent\textit{Wrapper, class and numbering.} The version of record is set with the Springer Nature class sn-jnl (options pdflatex, sn-mathphys-num), which the e-print ships as sn-jnl.cls together with the bibliography style sn-mathphys-num.bst. Neither is redistributed in this package, because the AMS wrapper of the pinned TeX Live runtime is used instead: the author\textquotesingle s own theorem environments (source lines 41--47), his two macros \textbackslash{}floor and \textbackslash{}at (lines 58--60) and his two enumitem lists mycases and mysubcases (lines 63--68) are carried over verbatim, while the class-specific \textbackslash{}theoremstyle\{thmstyleone\} line is replaced by the wrapper\textquotesingle s plain theorem style and the Springer class, its .bst, its author-year bibliography markup and the author\textquotesingle s algorithm/listings/table packages are not used. One editorial counter command restores the article\textquotesingle s own number of the context lemma (Lemma 5); the counted statement prints as Theorem 1, exactly as in the version PDF, since no other theorem environment precedes it. The article\textquotesingle s sequential equation numbering is not reproduced: the displays of this unit carry the wrapper\textquotesingle s own consecutive numbers, and every internal cross-reference resolves to the correct object inside the unit.\par
\noindent\textit{Citation map.} The delivered unit invokes no citation at all: neither the counted statement, nor its proof, nor the retained context lemma contains a \textbackslash{}cite command. The article\textquotesingle s bibliography is an external springer .bbl whose entries are written in the sn-jnl author-year macros (barticle, bauthor, batitle, ...) that only the non-redistributed class defines; since the two statements that do carry attributions (the imported mate lemma and the underlying GCP lemma) lie outside this unit, no bibliography entry is needed and the delivered bibliography is empty by construction.\par
\noindent\textit{Assets.} The e-print of this version ships the author .tex, an external .bbl, a .bib, the Springer class and .bst, and six .eps figures. None of them is redistributed. No retained span contains any \textbackslash{}includegraphics, \textbackslash{}input or \textbackslash{}include command, and no figure is delivered; the paper\textquotesingle s six figures lie outside the counted unit. Every mathematical symbol, sentence and label of the delivered spans is the author\textquotesingle s own native text; no mathematics was written, repaired, re-derived or normalized, and every printed slip inside these spans - including the author\textquotesingle s own wording such as the phrase the fact that 0 < tau\_1 \textbackslash{}lneqq 2\^{}\{m+1\} - is kept literal.\par
\setcounter{lemma}{4}
	\begin{lemma}\label{Thm_IGC_gen}
		The code set $\mathcal{S} = \{ \mathcal{C}_{\mathbf{s}, \mathbf{t}}: \mathbf{s}, \mathbf{t} \in \prod_{\alpha=1}^{k} \mathbb{Z}_{p_\alpha}\}$  forms a $(K,M, L,Z)$-IGC code set with $K= (p_1 p_2 \dots p_k)^2$, $M= p_1 p_2 \dots p_k$, $L= p_1^{m_1} p_2^{m_2} \dots p_k^{m_k}$ and $Z= p_1^{m_1-1} p_2^{m_2-1} \dots p_k^{m_k-1}$. The $\mathbf{t}$-th IGC code group is expressed as $I^{\mathbf{t}} = \{ \mathcal{C}_{\mathbf{s}, \mathbf{t}}: \mathbf{s} \in \prod_{\alpha=1}^{k} \mathbb{Z}_{p_\alpha}\}$.
	\end{lemma}	

	\begin{theorem}\label{thm_ZCACS}
		Each set $\mathcal{X}^{\bm{\zeta}}= \{\mathbf{X}^{\bm{\gamma},\bm{\zeta}}: \bm{\gamma} \in \prod_{\alpha=1}^{k} \mathbb{Z}_{p_\alpha}\}$ is a 2-D ZCAC of size  $L_1 \times L_2$, where $L_1= 2^m p$ and $L_2= 2 p_1^{m_1} p_2^{m_2} \dots p_k^{m_k}$, ZCZ size is $Z_1 \times Z_2$,  where $Z_1= 2^{m+1} $ and $Z_2=  p_1^{m_1-1} p_2^{m_2-1} \dots p_k^{m_k-1}$ and flock size $M= p_1 p_2 \dots p_k$. The set $\{ \mathcal{X}^{\bm{\zeta}}: \bm{\zeta} \in \Gamma\}$ is a 2-D ZCACS with set size $K= 2p_1^2 p_2^2 \dots p_k^2$.
	\end{theorem} 

	\begin{proof}
		To prove the theorem, we need to show that
		\begin{equation}
			\begin{split}
				\sum_{\bm{\gamma} \in \prod_{\alpha=1}^{k} \mathbb{Z}_{p_\alpha}}^{} C\left( \mathbf{X}^{\bm{\gamma},\bm{\zeta}_1}, \mathbf{X}^{\bm{\gamma},\bm{\zeta}_2} \right) (\tau_1, \tau_2)
				=
				\begin{cases}
					L_1 L_2 \prod_{\alpha=1}^{k} p_\alpha,& \bm{\zeta}_1 = \bm{\zeta}_2, (\tau_1, \tau_2 ) = (0,0);\\
					0, & \bm{\zeta}_1 = \bm{\zeta}_2, (\tau_1, \tau_2 ) \neq (0,0);\\
					& \left| \tau_1 \right| < 2^{m+1}, \\
					& \left| \tau_2 \right| < \prod_{\alpha=1}^{k} p_\alpha^{m_\alpha-1};\\
					0, & \bm{\zeta}_1 \neq \bm{\zeta}_2, \left| \tau_1 \right| < 2^{m+1},\\
					& \left| \tau_2 \right| < \prod_{\alpha=1}^{k} p_\alpha^{m_\alpha-1},\\
				\end{cases}
			\end{split}
		\end{equation}
		where $L_1=2^{m}p$ $ L_2= 2\prod_{\alpha=1}^{k} p_\alpha^{m_\alpha}$.
		It is straightforward to verify that the number of $\bm{\zeta} \in \Gamma $ is $K= 2 p_1^2 p_2^2 \dots p_k^2$. Also, for the trivial case $\bm{\zeta}_1 = \bm{\zeta}_2, (\tau_1, \tau_2 ) = (0,0)$, we have
		\begin{equation}
			\begin{split}
				\sum_{\bm{\gamma} \in \prod_{\alpha=1}^{k} \mathbb{Z}_{p_\alpha}}^{} C\left( \mathbf{X}^{\bm{\gamma},\bm{\zeta}_1}, \mathbf{X}^{\bm{\gamma},\bm{\zeta}_2} \right) (\tau_1, \tau_2)=	2^{m+1} p \prod_{\alpha=1}^{k} p_\alpha^{m_\alpha+1}.
			\end{split}
		\end{equation}
		To prove the rest of the theorem, we shall consider only $\tau_1 \geq 0$, and $0 \leq \left| \tau_2\right| < \prod_{\alpha=1}^{k} p_\alpha^{m_\alpha-1}$, because for any two 2-D arrays $\mathbf{X}_1$, $\mathbf{X}_2$, we have $C(\mathbf{X}_1, \mathbf{X}_2) (-\tau_1, \tau_2)= C(\mathbf{X}_2, \mathbf{X}_1)^* (\tau_1, -\tau_2)$. For a fixed $\bm{\gamma}$, we have
		\begin{equation}
			\begin{split}
				C\left( \mathbf{X}^{\bm{\gamma},\bm{\zeta}_1}, \mathbf{X}^{\bm{\gamma}, \bm{\zeta}_2} \right) (\tau_1, \tau_2)
				= &\sum_{i=0}^{L_1 -\tau_1-1} \sum_{j=0}^{L_2-\tau_2-1} \mathbf{X}^{\bm{\gamma}, \bm{\zeta}_1}_{i,j} {\mathbf{X}^{\bm{\gamma}, \bm{\zeta}_2}_{i', j'}}^*\\
				=& \sum_{i=0}^{L_1 -\tau_1-1} \sum_{j=0}^{L_2-\tau_2-1} \omega_q^{ 	F_{ \vec{i}}^{\bm{\gamma}, \bm{\zeta}_1} (\vec{j}) } \omega_q^{ 	-F_{ \vec{i'}}^{\bm{\gamma}, \bm{\zeta}_2} (\vec{j'}) },
			\end{split}
		\end{equation}
		where $i'=i+ \tau_1$, $j'=j+\tau_2$. Let $1+ \sum_{\alpha=1}^{k} m_\alpha= \epsilon$ and $m+1= \vartheta$. 
		We use the following notations:
		\begin{enumerate}
			\item We let $\vec{i}= (\vec{i}_{\delta}, i_{\vartheta}) $ and  $\vec{i'}= (\vec{i}'_{\delta}, i'_{\vartheta}) $, where $\vec{i}_{\delta}, \vec{i}'_{\delta}$ are the vectors with values in $\prod_{\alpha=1}^{k} \mathbb{Z}_{p_\alpha}^{m_\alpha-1} \times \prod_{\alpha=1}^{k} \mathbb{Z}_{p_\alpha}$ which are truncated from $\vec{i}$ and $\vec{i'}$, respectively.
			Whereas, $i_{\vartheta}, i'_{\vartheta}$ are the last elements of $\vec{i}, \vec{i}'$ corresponding to $\mathbb{Z}_2$.
			
			\item Similarly, we define $\vec{j}=(\vec{j}_{\xi}, j_{\epsilon})$ and $\vec{j}'=(\vec{j}'_{\xi}, j'_{\epsilon})$, where $\vec{j}_{\xi}, \vec{j}'_{\xi} \in \mathbb{Z}_2^m$ and $\vec{j}_{\epsilon}, \vec{j}'_{\epsilon} \in \mathbb{Z}_p$.  
			
			\item We define $i_{\delta}, i'_{\delta}$ to be the decimal numbers derived from the radix-representations $\vec{i}_{\delta}, \vec{i}'_{\delta}$, respectively. Similarly, statement holds for $j_{\xi}, j'_{\xi}$.
		\end{enumerate}
		Then, expanding the function $ F_{ \vec{d}'}^{\bm{\gamma}, \bm{\zeta}} (\vec{x})$, we have
		\begin{equation}\label{Thm_4_omega_eqn_1}
			\begin{split}
				\omega_q^{ F_{ \vec{i}}^{\bm{\gamma}, \bm{\zeta}_1} (\vec{j}) } \omega_q^{ -F_{ \vec{i}'}^{\bm{\gamma}, \bm{\zeta}_2} (\vec{j}') }
				=& \omega_q^{ (1-j_{\epsilon}) \left\{a_{\mathbf{s}_1^1,\mathbf{t}_1^1}^{\bm{\gamma}} (\vec{j}_{\xi}) + a_1(\vec{i}_{\delta}) W_1(i_\vartheta) + a_2(\vec{i}_{\delta}) W_2(i_\vartheta) \right\}} \\
				&  \times \omega_q^{ j_{\epsilon} \left\{a_{\mathbf{s}_2^1,\mathbf{t}_2^1}^{\bm{\gamma}} (\vec{j}_{\xi}) + b_1(\vec{i}_{\delta}) W_1(i_\vartheta) + b_2(\vec{i}_{\delta}) W_2(i_\vartheta)\right\}} \\
				& \times \omega_q^{ -(1-j'_{\epsilon}) \left\{a_{\mathbf{s}_1^2,\mathbf{t}_1^2}^{\bm{\gamma}} (\vec{j}'_{\xi}) + a_1(\vec{i}'_{\delta}) W_1(i'_\vartheta) + a_2(\vec{i}'_{\delta}) W_2(i'_\vartheta)\right\}} \\
				& \times \omega_q^{ -j'_{\epsilon} \left\{a_{\mathbf{s}_2^2,\mathbf{t}_2^2}^{\bm{\gamma}} (\vec{j}'_{\xi}) + b_1(\vec{i}'_{\delta}) W_1(i'_\vartheta) + b_2(\vec{i}'_{\delta}) W_2(i'_\vartheta)\right\}},
			\end{split}
		\end{equation}
		Now, we split the rest of the proof into three cases.
		
		\begin{mycases}
			%-----------------------case-1------------------
			\item Let  $0 \leq \tau_1 < 2^{m+1}$ and $0 < \tau_2 < \prod_{\alpha=1}^{k} p_\alpha^{m_\alpha-1}$. Then define
			\begin{equation}
				\begin{split}
					S_1=&  \sum_{j=0}^{\prod_{\alpha=1}^{k} p_\alpha^{m_\alpha} -\tau_2-1} \omega_q^{ 	F_{ \vec{i}}^{\bm{\gamma}, \bm{\zeta}_1} (\vec{j}) } \omega_q^{ 	-F_{ \vec{i}'}^{\bm{\gamma}, \bm{\zeta}_2} (\vec{j}') },\\
					S_2=& \sum_{j= \prod_{\alpha=1}^{k} p_\alpha^{m_\alpha} -\tau_2 }^{\prod_{\alpha=1}^{k}  p_\alpha^{m_\alpha}-1} \omega_q^{ 	F_{ \vec{i}}^{\bm{\gamma}, \bm{\zeta}_1} (\vec{j}) } \omega_q^{ 	-F_{ \vec{i}'}^{\bm{\gamma}, \bm{\zeta}_2} (\vec{j}') },\\
					S_3=& \sum_{j= \prod_{\alpha=1}^{k} p_\alpha^{m_\alpha}}^{2 \prod_{\alpha=1}^{k} p_\alpha^{m_\alpha} -\tau_2-1} \omega_q^{ 	F_{ \vec{i}}^{\bm{\gamma}, \bm{\zeta}_1} (\vec{j}) } \omega_q^{ 	-F_{ \vec{i}'}^{\bm{\gamma}, \bm{\zeta}_2} (\vec{j}') }.\\
				\end{split}
			\end{equation}
			Then, we can write the term $ C\left( \mathbf{X}^{\bm{\gamma},\bm{\zeta}_1}, \mathbf{X}^{\bm{\gamma}, \bm{\zeta}_2} \right) (\tau_1, \tau_2)$ in the following form:
			\begin{equation}\label{eqn_S_sum}
				C\left( \mathbf{X}^{\bm{\gamma},\bm{\zeta}_1}, \mathbf{X}^{\bm{\gamma}, \bm{\zeta}_2} \right) (\tau_1, \tau_2)= \sum_{i=0}^{L_1 -\tau_1-1} \bigg[ S_1 + S_2+ S_3 \bigg].
			\end{equation}
			
			%---------------------S-1----------------
			\begin{enumerate}
				\item For $S_1$, we have $0 \leq j \leq \prod_{\alpha=1}^{k} p_\alpha^{m_\alpha} -\tau_2-1$. This implies that $j_{\epsilon}=j'_{\epsilon}=0$ as $\tau_2 < \prod_{\alpha=1}^{k} p_\alpha^{m_\alpha-1}$. Similarly, we can find that  $0 \leq j_{\xi} \leq \prod_{\alpha=1}^{k} p_\alpha^{m_\alpha} -\tau_2-1$,  $\tau_2 \leq j'_{\xi} \leq \prod_{\alpha=1}^{k} p_\alpha^{m_\alpha}-1$. So, using (\ref{Thm_4_omega_eqn_1}), we get
				\begin{equation}
					\begin{split}
						S_1= &\omega_q^{a_1(\vec{i}_{\delta}) W_1(i_\vartheta) + a_2(\vec{i}_{\delta}) W_2(i_\vartheta)} \omega_q^{- \left( a_1(\vec{i}'_{\delta}) W_1(i'_\vartheta) + a_2(\vec{i}'_{\delta}) W_2(i'_\vartheta) \right)}\\
						& \times \left( \sum_{j_{\xi}=0}^{\prod_{\alpha=1}^{k} p_\alpha^{m_\alpha} -\tau_2-1} \omega_q^{a_{\mathbf{s}_1^1,\mathbf{t}_1^1}^{\bm{\gamma}} (\vec{j}_{\xi})-a_{\mathbf{s}_1^2,\mathbf{t}_1^2}^{\bm{\gamma}} (\vec{j}'_{\xi}) } \right)\\
					\end{split}
				\end{equation}
				Then, whether or not $\mathbf{s}_1^1= \mathbf{s}_1^2$ or $\mathbf{t}_1^1= \mathbf{t}_1^2$, i.e., irrespective of $\bm{\zeta}_1 = \bm{\zeta}_2$ or $\bm{\zeta}_1 \neq \bm{\zeta}_2$, from \textit{Lemma \ref{Thm_IGC_gen}}, we can conclude
				\begin{equation}\label{eqn_S1}
					\sum_{\bm{\gamma} \in \prod_{\alpha=1}^{k} \mathbb{Z}_{p_\alpha}}^{} \sum_{j_{\xi}=0}^{\prod_{\alpha=1}^{k} p_\alpha^{m_\alpha} -\tau_2-1} \omega_q^{a_{\mathbf{s}_1^1,\mathbf{t}_1^1}^{\bm{\gamma}} (\vec{j}_{\xi})-a_{\mathbf{s}_1^2,\mathbf{t}_1^2}^{\bm{\gamma}} (\vec{j}'_{\xi}) }=0,
				\end{equation}
				because $0 \lneqq \tau_2 <  \prod_{\alpha=1}^{k} p_\alpha^{m_\alpha-1}$.
				
				%----------S-2---------------
				\item For $S_2$, using a similar argument, we get
				\begin{equation}\label{eqn_S2}
					\begin{split}
						\sum_{\bm{\gamma} \in \prod_{\alpha=1}^{k} \mathbb{Z}_{p_\alpha}}^{} \sum_{j_{\xi}=\prod_{\alpha=1}^{k} p_\alpha^{m_\alpha} -\tau_2}^{\prod_{\alpha=1}^{k}  p_\alpha^{m_\alpha}-1} \omega_q^{a_{\mathbf{s}_1^1,\mathbf{t}_1^1}^{\bm{\gamma}} (\vec{j}_{\xi})-a_{\mathbf{s}_2^2,\mathbf{t}_2^2}^{\bm{\gamma}} (\vec{j}'_{\xi}) }=0.
					\end{split}
				\end{equation}
				
				%-----------------S-3--------------------
				\item For $S_3$, similarly we have
				\begin{equation}\label{eqn_S3}
					\begin{split}
						\sum_{\bm{\gamma} \in \prod_{\alpha=1}^{k} \mathbb{Z}_{p_\alpha}}^{} \sum_{j_{\xi}= \prod_{\alpha=1}^{k} p_\alpha^{m_\alpha}}^{2 \prod_{\alpha=1}^{k} p_\alpha^{m_\alpha} -\tau_2-1} \omega_q^{ 	F_{ \vec{i}}^{\bm{\gamma}, \bm{\zeta}_1} (\vec{j}_{\xi}) } \omega_q^{ 	-F_{ \vec{i}'}^{\bm{\gamma}, \bm{\zeta}_2} (\vec{j}'_{\xi}) }=0.
					\end{split}
				\end{equation}
			\end{enumerate}
			Now, combining (\ref{eqn_S1}), (\ref{eqn_S2}), (\ref{eqn_S3}), and (\ref{eqn_S_sum}), we have
			\begin{equation}
				\begin{split}
					\sum_{\bm{\gamma} \in \prod_{\alpha=1}^{k} \mathbb{Z}_{p_\alpha}}^{} C\left( \mathbf{X}^{\bm{\gamma},\bm{\zeta}_1}, \mathbf{X}^{\bm{\gamma},\bm{\zeta}_2} \right) (\tau_1, \tau_2)=0.
				\end{split}
			\end{equation}
			
			%================ Case-2=====================
			\item  Let $\tau_2=0$, and $0 < \tau_1 < 2^{m+1}$. Then, we have
			\begin{equation}\label{eqn_T_sum}
				\begin{split}
					C\left( \mathbf{X}^{\bm{\gamma},\bm{\zeta}_1}, \mathbf{X}^{\bm{\gamma}, \bm{\zeta}_2} \right) (\tau_1, \tau_2) =  \sum_{i=0}^{L_1 -\tau_1-1} \bigg[ T_1 + T_2 \bigg],\\
				\end{split}
			\end{equation}
			where
			\begin{equation}
				\begin{split}
					T_1= & \sum_{j=0}^{\prod_{\alpha=1}^{k} p_\alpha^{m_\alpha} -1} \omega_q^{ 	F_{ \vec{i}}^{\bm{\gamma}, \bm{\zeta}_1} (\vec{j}) } \omega_q^{ 	-F_{ \vec{i}'}^{\bm{\gamma}, \bm{\zeta}_2} (\vec{j}) },\\
					T_2= & \sum_{j= \prod_{\alpha=1}^{k} p_\alpha^{m_\alpha}}^{2 \prod_{\alpha=1}^{k} p_\alpha^{m_\alpha} -1} \omega_q^{ 	F_{ \vec{i}}^{\bm{\gamma}, \bm{\zeta}_1} (\vec{j}) } \omega_q^{ 	-F_{ \vec{i}'}^{\bm{\gamma}, \bm{\zeta}_2} (\vec{j}) }.
				\end{split}
			\end{equation}
			As $\tau_2=0$, this implies that for $T_1$ we have $j_{\epsilon}=0$, and for $T_2$ we have $j_{\epsilon}=1$.
			Hence, from (\ref{Thm_4_omega_eqn_1}), we have
			\begin{equation}\label{eqn_W1W2_1}
				\begin{split}
					T_1= &\omega_q^{a_1(\vec{i}_{\delta}) W_1(i_\vartheta) + a_2(\vec{i}_{\delta}) W_2(i_\vartheta)} \omega_q^{- \left( a_1(\vec{i}'_{\delta}) W_1(i'_\vartheta) + a_2(\vec{i}'_{\delta}) W_2(i'_\vartheta) \right)}\\
					& \times \left( \sum_{j_{\xi}=0}^{\prod_{\alpha=1}^{k} p_\alpha^{m_\alpha} -1} \omega_q^{a_{\mathbf{s}_1^1,\mathbf{t}_1^1}^{\bm{\gamma}} (\vec{j}_{\xi})-a_{\mathbf{s}_1^2,\mathbf{t}_1^2}^{\bm{\gamma}} (\vec{j}_{\xi}) } \right),\\
				\end{split}
			\end{equation}
			and
			\begin{equation}\label{eqn_W1W2_2}
				\begin{split}
					T_2=& \omega_q^{b_1(\vec{i}_{\delta}) W_1(i_\vartheta) + b_2(\vec{i}_{\delta}) W_2(i_\vartheta)}  \omega_q^{- \left( b_1(\vec{i}'_{\delta}) W_1(i'_\vartheta) + b_2(\vec{i}'_{\delta}) W_2(i'_\vartheta) \right)}\\
					& \times \left( \sum_{j_{\xi}= \prod_{\alpha=1}^{k} p_\alpha^{m_\alpha}}^{2 \prod_{\alpha=1}^{k} p_\alpha^{m_\alpha} -1} \omega_q^{a_{\mathbf{s}_2^1,\mathbf{t}_2^1}^{\bm{\gamma}} (\vec{j}_{\xi})-a_{\mathbf{s}_2^2,\mathbf{t}_2^2}^{\bm{\gamma}} (\vec{j}_{\xi}) } \right).\\
				\end{split}
			\end{equation}
			Now, we have two sub-cases.
			
			\begin{mysubcases}
				\item Let $\bm{\zeta}_1 = \bm{\zeta}_2$. In that case, we have
				\begin{equation}
					\begin{split}
						\sum_{j_{\xi}=0}^{\prod_{\alpha=1}^{k} p_\alpha^{m_\alpha} -1} \omega_q^{a_{\mathbf{s}_1^1,\mathbf{t}_1^1}^{\bm{\gamma}} (\vec{j}_{\xi})-a_{\mathbf{s}_1^2,\mathbf{t}_1^2}^{\bm{\gamma}} (\vec{j}_{\xi}) }=& \prod_{\alpha=1}^{k} p_\alpha^{m_\alpha},\\
						\sum_{j_{\xi}= \prod_{\alpha=1}^{k} p_\alpha^{m_\alpha}}^{2 \prod_{\alpha=1}^{k} p_\alpha^{m_\alpha} -1} \omega_q^{a_{\mathbf{s}_2^1,\mathbf{t}_2^1}^{\bm{\gamma}} (\vec{j}_{\xi})-a_{\mathbf{s}_2^2,\mathbf{t}_2^2}^{\bm{\gamma}} (\vec{j}_{\xi}) }=& \prod_{\alpha=1}^{k} p_\alpha^{m_\alpha}.\\
					\end{split}
				\end{equation}
				Also, using the property of $W_1$ and $W_2$, we can find that
				\begin{equation}
					\begin{split}
						W_1(i_\vartheta)=0  \Leftrightarrow  &~ W_2(i_\vartheta)=1,\\
						W_1(i_\vartheta)=1 \Leftrightarrow  &~ W_2(i_\vartheta)=0.\\
					\end{split}
				\end{equation}
				Hence, without loss of generality, from the fact that $0< \tau_1 \lneqq 2^{m+1}$ and from (\ref{eqn_W1W2_1}) and (\ref{eqn_W1W2_2}), we get
				\begin{equation}
					\begin{split}
						\sum_{i=0}^{L_1 -\tau_1-1} \bigg[ T_1 + T_2 \bigg]
						=&\prod_{\alpha=1}^{k} p_\alpha^{m_\alpha} \sum_{i_{\delta}=0}^{L_1 -\tau_1-1} \bigg[ \left( \omega_q^{ a_1 (\vec{i}_{\delta})- a_1 (\vec{i}'_{\delta})} + \omega_q^{ b_1 (\vec{i}_{\delta})- b_1 (\vec{i}'_{\delta})}  \right)\\
						& + \left( \omega_q^{ a_2 (\vec{i}_{\delta})- a_1 (\vec{i}'_{\delta})} + \omega_q^{ b_2 (\vec{i}_{\delta})- b_1 (\vec{i}'_{\delta})}  \right)\\ &+\left( \omega_q^{ a_2 (\vec{i}_{\delta})- a_2 (\vec{i}'_{\delta})} + \omega_q^{ b_2 (\vec{i}_{\delta})- b_2 (\vec{i}'_{\delta})}  \right)   \bigg]\\
						&=0,
					\end{split}
				\end{equation}
				using Lemma \href{https://arxiv.org/abs/2109.00970v4}{9}. Therefore
				\begin{equation}
					C\left( \mathbf{X}^{\bm{\gamma},\bm{\zeta}_1}, \mathbf{X}^{\bm{\gamma}, \bm{\zeta}_2} \right) (\tau_1, \tau_2)=0.
				\end{equation}
				
				\item Let $\bm{\zeta}_1 \neq \bm{\zeta}_2$. If $(\mathbf{s}_i^1, \mathbf{t}_i^1)= (\mathbf{s}_i^2, \mathbf{t}_i^2)$, then we  proceed similar to the previous case and get
				\begin{equation}
					\begin{split}
						C\left( \mathbf{X}^{\bm{\gamma},\bm{\zeta}_1}, \mathbf{X}^{\bm{\gamma}, \bm{\zeta}_2} \right) (\tau_1, \tau_2)=0.
					\end{split}
				\end{equation}
				However, if either $\mathbf{s}_i^1 \neq \mathbf{s}_i^2$ or $\mathbf{t}_i^1 \neq \mathbf{t}_i^2$, then using Lemma \ref{Thm_IGC_gen}, we have
				\begin{equation}
					\begin{split}
						\sum_{\bm{\gamma} \prod_{\alpha=1}^{k} \mathbb{Z}_{p_\alpha}}^{ }	\sum_{j_{\xi}}^{} \omega_q^{a_{\mathbf{s}_i^1,\mathbf{t}_i^1}^{\bm{\gamma}} (\vec{j}_{\xi})-a_{\mathbf{s}_i^2,\mathbf{t}_i^2}^{\bm{\gamma}} (\vec{j}_{\xi}) }=0
					\end{split}.
				\end{equation} 
			\end{mysubcases}
			
			%--------------------------case-2------------
			\item Let  $0 \leq \tau_1 < 2^{m+1}$ and $ -\prod_{\alpha=1}^{k} p_\alpha^{m_\alpha-1} < \tau_2 < 0$. This case can be proved similar to \textit{Case I}.
		\end{mycases}
		Combining all these cases, we have the theorem.
	\end{proof}

\end{document}
